\documentclass[a4paper,10pt,conference]{ieeeconf}
\usepackage[utf8]{inputenc}
\usepackage{graphicx}
\usepackage{amsmath,amssymb}
\usepackage{url}
\usepackage{xcolor}
\usepackage{hyperref}
\hypersetup{colorlinks=true,urlcolor=blue,citecolor=blue,linkcolor=black}
\urlstyle{same}
% Organizer header retained.
\title{WACV 2027 UPAR Challenge\\ \textit{Fact sheet}}
\begin{document}
\maketitle
This is the fact sheet's template for the \href{https://vap.aau.dk/rws/challenge/}{WACV 2027 UPAR Challenge}. Please fill out the following sections carefully in a scientific writing style. Then, send the compressed project (in \texttt{.zip} format), i.e., the generated PDF, \texttt{.tex}, \texttt{.bib} and any additional files to \url{upar.challenge@gmail.com}.

\section{Team details}\label{teamdetails}
\begin{itemize}
\item Team leader name: Kaveh Safavigerdini.
\item Username(s) on Codabench: \\ \texttt{kavehsfv-994308-994308\_VID-Docker-GPU}.
\item Team leader affiliation: Electrical Engineering and Computer Science Department, University of Missouri--Columbia, USA.
\item Team leader email: \url{ksgh2@umsystem.edu}.
\item Name of other team members (and affiliation): None.
\item Team website URL (if any): \url{https://github.com/kavehsfv}.
\item Competition track (mark with X one single option)\footnote{If you participated in more than one competition track, you need to share with the organizers one fact sheet per track.}:
\begin{itemize}
\item ( ) Track 1: Pedestrian Attribute Recognition.
\item ( ) Track 2: Attribute-based Person Retrieval.
\item (X) Track 3: Human Pose Estimation.
\end{itemize}
\end{itemize}

\section{Contribution details}
\subsection{Title of the contribution}
\textbf{Expected-OKS pose decoding with ViTPose+ Small, reduced-dose box jitter and multi-scale flip inference.}
The final entry (\texttt{mipnet\_v3}) fine-tunes a pretrained top-down pose estimator with probabilistic heatmaps, then combines three crop scales. It uses a single set of model weights. The name denotes our experiment identifier; it is not the original multi-instance MIPNet architecture. This fact sheet describes only the submitted final method.

\subsection{Representative image / workflow diagram of the method}
\begin{figure}[ht]
\centering
\setlength{\fboxsep}{5pt}
\fbox{\begin{minipage}{0.91\columnwidth}\centering
Image + organizer-provided person box\\
$\downarrow$\\
Aspect correction + padding $1.25$\\
Three crop scales $\{0.9,1.0,1.1\}$\\
RGB $256\!\times\!192$ + horizontal flips\\
$\downarrow$\\
Shared ViTPose+ Small (COCO route 0)\\
PT18 remap/refine $\rightarrow$ Sparsemax heatmaps\\
Detached branch $\rightarrow$ in-crop presence\\
$\downarrow$\\
Undo flips; average paired heatmaps\\
OKS-kernel convolution + Taylor refinement\\
$\downarrow$\\
Map each scale to image coordinates\\
Average coordinates, confidence and presence\\
17 $(x,y,v)$ joints + one score per box
\end{minipage}}
\caption{Final inference pipeline. Six views are stacked into one model call per box chunk; the three coordinate estimates are fused only after mapping to image space.}
\label{fig:pipeline}
\end{figure}

\subsection{Detailed method description}
\paragraph{Model and representation.}
We use the starter's public \texttt{usyd-community/vitpose-plus-small} checkpoint~\cite{startermodel}, based on ViTPose~\cite{xu2022vitpose}. Expert index 0 is supplied for every crop. The pretrained decoder produces 17 heatmaps at $64\!\times\!48$. A learned $1\!\times\!1$ channel projection, residual depthwise $3\!\times\!3$ refinement and final $1\!\times\!1$ projection adapt COCO channels to the 17 PoseTrack18 output slots. An additional $1\!\times\!1$ refinement followed by spatial Sparsemax with temperature $0.5$ yields nonnegative unit-mass maps $P_k$. The remapping is initialized from anatomical COCO-to-PT18 correspondences, with head landmarks initialized using available face/shoulder proxies.

Following adapted ProbPose mechanisms~\cite{purkrabek2025probpose}, a separate branch predicts in-crop presence $r_k$. Its input heatmap features are detached, so its BCE gradient does not change the localization pathway. Three $3\!\times\!3$ convolution/BN/pooling/ReLU stages, a $1\!\times\!1$ convolution, sigmoid and spatial averaging produce 17 presence values. Presence is not a calibrated localization-error estimate.

\paragraph{Preprocessing and data.}
Boxes are supplied by the organizer; no detector is trained or run. Each box is adjusted to aspect ratio $192/256$, expanded by $1.25$ and sampled by a bilinear affine warp into an RGB $256\!\times\!192$ crop. Normalization uses mean $(0.485,0.456,0.406)$ and standard deviation $(0.229,0.224,0.225)$.

We fine-tune on the supplied \texttt{train\_upar\_all.json} snapshot containing 42,736 annotations. The included 301-annotation local split is excluded from gradient updates, leaving 42,435 training annotations: 29,613 PA-100K and 12,822 Market. No PETA annotation occurs in this snapshot. No official validation or hidden-test image is added to training. The local holdout selects the checkpoint by crop-space pixel error; it is not an additional training set.

\paragraph{Targets and losses.}
Let $u_k$ be a labeled joint in heatmap coordinates. We use
\begin{align}
 A_h &= \sqrt{(64/1.25)(48/1.25)},\nonumber\\
 s_k &= \operatorname{clip}\!\left(2(2\sigma_k)^2 A_h,\,0.55,\,3.0\right),\\
 T_k(p) &= \exp\!\left[-\|p-u_k\|^2/(2s_k)\right].
\end{align}
Here $s_k$ is a variance-like denominator parameter, not a standard deviation; $A_h\approx44.3405$ is the fixed crop-derived quantity used in our implementation. In output order the sigmas are
\begin{equation*}
\begin{split}
( & .025,.025,.059,.059,.059,.079,.079,.072,.072,\\
 & .062,.062,.107,.107,.087,.087,.089,.089).
\end{split}
\end{equation*}
Targets are not normalized to unit mass, whereas predictions are.

With $w_k$ the annotation mask, the localization loss is
\begin{equation}
\begin{split}
 L_{\mathrm{loc}} &= \frac{1}{64\cdot48\sum_k w_k}\sum_{k,p} w_k\bigl\{\\
 &0.95P_k(p)[1-T_k(p)]+0.05\|\nabla_S P_k(p)\|^2\bigr\}.
\end{split}
\end{equation}
where $\nabla_S$ uses horizontal and vertical Sobel filters. The overall objective is $L=L_{\mathrm{loc}}+0.1L_{\mathrm{BCE}}(r,z)$, with $z_k$ indicating whether the transformed joint lies inside the crop and BCE masked by $w_k$. Unlabeled joints contribute neither loss. This is an adapted probabilistic heatmap loss, not the complete original ProbPose system.

\paragraph{Augmentation and optimization.}
With probability $0.25$, the box center is shifted uniformly within $\pm5\%$ of its width/height and each side length is perturbed within $\pm7.5\%$. Source-image keypoints remain unchanged; only the crop window moves. Additional augmentation comprises flip probability $0.5$, Gaussian scale $\mathcal{N}(1,0.1^2)$ clipped to $[0.7,1.3]$, uniform rotation $\pm30^\circ$, brightness/contrast probability $0.5$ (multiplicative $[0.8,1.2]$, additive $[-0.1,0.1]$), and one-sided blackout probability $0.3$.

Training lasts 15 epochs, batch size 64, seed 0, eight loader workers, AdamW weight decay $10^{-4}$ and BF16 autocast. During epochs 0--3 the pretrained network is frozen; the probability/refinement and remapping learning rates are $10^{-3}$ and $5\!\times\!10^{-4}$. At epoch 4 all pretrained parameters are unfrozen at $10^{-5}$, while the other two rates become $10^{-4}$ and $5\!\times\!10^{-5}$. No EMA or cosine schedule is used in the selected recipe. Presence BCE is computed in FP32 outside autocast. The selected checkpoint is epoch 13 (zero-based), with held-out crop-space error $3.452812$ pixels. This selection statistic is not official AP.

The historical training stack is Python 3.11, PyTorch 2.14.0+cu130, Transformers 5.18.0, NumPy 2.3.5 and Pillow 12.3.0. Training configuration, split, curve, checkpoint and offline export script are released. Inference is separately checked under PyTorch 2.4.1 CPU, the challenge-compatible generation.

\paragraph{Inference and output.}
At each scale, the original and flipped crops use the same weights. Flipped maps are spatially mirrored and left/right channels exchanged, without an additional pixel shift. Paired maps and presence estimates are averaged. Each map is convolved with a normalized Gaussian OKS kernel using the corresponding $s_k$; the maximum of this expected-OKS surface is refined with a second-order Taylor step. The resulting position is mapped from heatmap to crop to original-image coordinates. Coordinates, convolved-peak confidences and presence values are averaged across the three scales; coordinates are clamped to image bounds. The box score is the mean of its 17 averaged convolved peaks. We output $v_k=1$ for averaged presence $\geq0.5$, otherwise 0. Ear slots remain present although the challenge excludes them from scoring.

Up to eight boxes (48 views) are stacked per network call. Inference requires only PyTorch, NumPy and Pillow and performs no downloads. Degenerate boxes use an included training-pose prior. Multi-scale fusion takes place in image coordinates, not heatmap coordinates: an identical heatmap displacement represents a different image displacement at different crop scales.

\subsection{Challenge results}
\begin{table}[ht]
\centering
\caption{Official organizer leaderboard export, test set (36,509). Values are on the supplied $[0,1]$ scale.}
\begin{tabular}{lr}\hline
Metric & Final submission\\\hline
AP (mean) & \textbf{0.8862830000}\\
AP50 & 0.9672830000\\
AP\_medium & 0.8564020000\\
AP\_large & 0.6790190000\\\hline
\end{tabular}
\label{tab:results}
\end{table}
The participant confirms that the best/final method was uploaded; Table~\ref{tab:results} reports the official result of that entry. The release preserves its runtime-fixed inference payload. Checkpoint loading, both training stages, offline export parity and the inference contract have been verified locally. These checks are distinct from the organizer's hidden-set accuracy evaluation.

\subsection{Final remarks}
\textbf{Advantages:} one model, explicit PT18 adaptation, crop-noise training, metric-shaped decoding and self-contained offline inference. The detached presence branch avoids an auxiliary-gradient conflict with localization. \textbf{Limitations:} six test views increase compute; performance remains dependent on supplied boxes; convolved-peak scores are not calibrated error probabilities. Docker recipes are supplied but could not be built on the preparation host, where Docker is unavailable. Starter-checkpoint pretraining provenance requires the clarification below.

\section{Additional method details}
Please, reply if your challenge entry considered (or not) the following strategies and provide a brief explanation. For each question, mark with X one single option.
\begin{itemize}
\item \textbf{For the competition track associated with this fact sheet, you confirm that you have trained your model on the predefined datasets} (X) Yes, ( ) No.\\
Fine-tuning uses only the supplied training annotations, excluding the released 301-annotation local holdout.

\item \textbf{Did you use any pre-trained model:} (X) Yes, ( ) No.\\
The starter's \nolinkurl{usyd-community/vitpose-plus-small}, revision \texttt{0c30b6534bb6} (full revision in the release), with expert route 0. Its public model card describes multi-dataset pretraining including COCO, AI Challenger and MPII. Route 0 does not prove COCO-only pretraining, because shared parameters retain supervision from other datasets. The starter's stated COCO/PoseTrack18 restriction therefore needs organizer clarification for its own supplied checkpoint; we do not assert that broader pretraining has been removed.

\item \textbf{Did you use external data?} ( ) Yes, (X) No.\\
No additional external images or pseudo-labels were introduced during our fine-tuning. Broader upstream pretraining is separately disclosed above; this answer does not assert a COCO-only warm start.

\item \textbf{Did you perform any data augmentation?} (X) Yes, ( ) No.\\
Reduced-dose box jitter, scale, rotation, flip, color adjustment and one-sided blackout as specified in Section II.

\item \textbf{Did you use the provided validation set as part of your training set?} ( ) Yes, (X) No.\\
No official validation images enter gradient updates. The local training-data holdout selects the checkpoint only.

\item \textbf{Did you use any regularization strategies/terms?} (X) Yes, ( ) No.\\
AdamW weight decay $10^{-4}$ and Sobel heatmap smoothing weight $0.05$, in addition to augmentation.

\item \textbf{Did you use any pre-processing steps?} (X) Yes, ( ) No.\\
Aspect correction, padding, affine RGB resampling and ImageNet normalization.

\item \textbf{Did you use any post-processing steps?} (X) Yes, ( ) No.\\
Flip and scale fusion, OKS-kernel/Taylor decoding, inverse coordinate mapping and image-bound clamping.

\item \textbf{Did you use Vision Language Models (VLMs) of any kind?} ( ) Yes, (X) No.\\
The pipeline is vision-only; no language encoder or text supervision is used in our method.
\end{itemize}

\section{Code repository}
\textbf{Code repository:} \url{https://github.com/kavehsfv/UPAR_Track3}.\\
Inference and training code, exact configuration, split and trained state dictionary are supplied. Git LFS hosts both the training checkpoint and ready-to-run TorchScript weights. A separate weight ZIP and the compressed fact-sheet project are available at \url{https://github.com/kavehsfv/UPAR_Track3/releases/tag/v1.0}. The README gives pinned dependencies, training/export commands, offline inference validation and inference/training Dockerfiles. Deployment-payload checksums and local verification logs are included. Dataset images are not redistributed.

\bibliographystyle{IEEEtran}
\bibliography{references}
\end{document}
